\addchap{\lsAbbreviationsTitle}
% \addchap{Abbreviations and symbols}
Terms referring to specific morphosyntactic categories of the Teop language start with a capital letter.

The morphosyntactic functions and meanings of glossed morphemes, especially those that are not dealt with in Chapter 3 like the demonstratives (\textsc{dem}1 – \textsc{dem}7), can easily be searched for in the Multifunctional Teop-English Dictionary (\cite{Mosel:2019bt}), see \sectref{sec:mted} for filtered searches, and \url{https://dictionaria.clld.org/contributions/teop\#twords2}.\\

\begin{tabularx}{\textwidth}{lX}
    \textsc{acaus}&Anticausative Prefix, see \sectref{sec:morphology}\\
    \textsc{adjr}&The suffix -\textit{ana} derives Adjectives, see \sectref{sec:morphology}\\
    \textsc{advers}&Adversative Conjunction \textit{he} `but, while, when', see \sectref{sec:depclaus}\\
    \textsc{advr}&The prefix \textit{va}- derives Adverbs from Adjectives, see \sectref{sec:modjux} and \sectref{sec:morphology}\\
    \textsc{ana}&Anaphoric Demonstrative in Noun Phrases, see \sectref{sec:npintro}\\
    \textsc{appl}&Applicative Particle \textit{ni} in the VC is transitive, see \sectref{sec:phrasefuncclause}\\
    \textsc{art}&Article, see \sectref{sec:paraarticles} and \sectref{sec:classnphead}\\
    \textsc{caus}&Causative Prefix \textit{vaa}-, \textit{va}-, see \sectref{sec:morphology}\\
    \textsc{compl1}&Complementizer \textit{tea} introduces Complement Clauses with anaphoric zero subjects, see \sectref{sec:depclaus} \\
    \textsc{compl2}&Complementizer \textit{be}\textsuperscript{2} introduces Complement Clauses with a subject Noun Phrase or Pronoun that is different from that of the main clause, see \sectref{sec:depclaus}\\
    \textsc{consec}&Conjunction \textit{re} marks dependent non-embedded clauses that denote a consecutive event, see \sectref{sec:consecpron} and \sectref{sec:depclaus}\\
    \textsc{cont}&Adverb \textit{bata} expresses continuity; it is translated by `along, continuously, for a while' or left untranslated, see \sectref{sec:verbcomplintro}\\
    \textsc{dat}&Preposition \textit{ki}, \textit{k}-  indicates that its complement Pronoun or Noun Phrase refers to an addressee or beneficiary; it is often incorporated in the verb complex, see \sectref{sec:incorporated}\\
    \textsc{dem}&Demonstrative functioning as an argument, adjunct, or determiner; the distinct types of demonstrative are indicated by numbers from \textsc{dem}1 to \textsc{dem}7.\\
\end{tabularx}
\begin{tabularx}{\textwidth}{lX}
    \textsc{derel}&Derelational Suffix -\textit{na} derives non-relational Nouns from a relational Nouns, see \sectref{sec:posscon}\\
    \textsc{dim}&Diminutive Particle, see \sectref{sec:preheadnp}\\
    \textsc{dyad}&Dyadic Relationship Marker, see \sectref{sec:classnphead}\\
    \textsc{emph2}&The emphatic particle \textit{kou}\\
    \textsc{ex}&exclusive (`we, but not you'), see \sectref{sec:fourcat}\\
    \textsc{goal}&Preposition \textit{vo} governs an Adverb, a Locative Phrase or a Prepositional Phrase, see \sectref{sec:phrasefuncclause} and \sectref{sec:incorporated}\\
    \textsc{h}&the prefix \textit{e}- marks Heavy Personal Pronouns, see \sectref{sec:fourcat} and \sectref{sec:heavypron}\\
    \textsc{in}&inclusive (`we including you'), see \sectref{sec:fourcat}\\
    \textsc{ipfv}&Imperfective Aspect Marker, see \sectref{sec:aspekt}\\
    \textsc{indef}&Indefiniteness Marker in Noun Phrases, see \sectref{sec:preheadnp}\\
    \textsc{kin.plm}&Kinship Plural Marker, see \sectref{sec:classnphead}\\
    \textsc{lnk}&the Linker \textit{vaa} links modifying Adverbs, Locative Phrases, Prepositional Phrases and \textit{tea}-Complement Clauses to the modified Noun or Verb. Linker Phrases also function as the core of a Verb Complex, see \sectref{sec:subphrase}\\
    \textsc{mult}&the prefix \textit{va}- derives verbs denoting multiple situations, including reciprocal actions, see \sectref{sec:morphology}\\
    \textsc{nspec}&Non-Specific Article, see \sectref{sec:paraarticles}\\
    \textsc{neg}1 ... \textsc{neg}2&Double Negator \textit{sa(ka)} ... \textit{haa}, see \sectref{sec:subphrase} and \sectref{sec:verbcomplintro}\\
    \textsc{obj.art}&Object Article introduces non-topical object NPs, see \sectref{sec:verbal} and \sectref{sec:paraarticles}\\
    \textsc{objm}&Object Marker within the Verb Complex, see \sectref{sec:verbcomplintro}\\
    \textsc{ord}&Ordinal prefix \textit{vaa}-, see \sectref{sec:morphology}\\
    \textsc{part}&Partitive (article), see \sectref{sec:paraarticles}\\
    \textsc{pfv}&Perfective Aspect Marker within the Verb Complex, see \sectref{sec:aspekt}\\
    \textsc{pl}&plural, see \sectref{sec:verbcomplintro}, \sectref{sec:paraarticles}, \sectref{sec:fourcat}, \tabref{tab:thirdpronfunc}, and \tabref{tab:fourthperson}\\
    \textsc{plm}&Plural Marker heading NPs, see \sectref{sec:classnphead}\\
    \textsc{poss}&Possessive Marker in Inalienable Possessive Constructions, see \sectref{sec:fourcat} and \sectref{sec:posscon}\\
    \textsc{poss.pron}&Possessive Pronoun, see \sectref{sec:fourcat} and \sectref{sec:posscon}\\
    \textsc{prep}&the multi-purpose Preposition \textit{te}, see \sectref{sec:subphrase}\\
    \textsc{pron}&Pronoun, see \tabref{tab:thirdpronfunc} in \sectref{sec:fourcat}\\
\end{tabularx}
\begin{tabularx}{\textwidth}{lX}
    \textsc{red}&reduplication, see \sectref{sec:morphology}\\
    \textsc{rel}&Relative Pronoun, see \sectref{sec:depclaus}\\
    \textsc{sg}&singular, see \sectref{sec:npintro}, \sectref{sec:fourcat}, \tabref{tab:thirdpronfunc}, and \tabref{tab:fourthperson}\\
    \textsc{tm}&the pre-head particles \textsc{tm}1--\textsc{tm}8 in Verb Complexes mark temporal relations and modality, see \sectref{sec:tmmarker}\\
\end{tabularx}